% =====================================================================
%  TEMPLATE -- copy to lectures/lecNN-short-name.tex and add
%  \include{lectures/lecNN-short-name} to main.tex.
%  (This file itself is NOT included in main.tex.)
%
%  STYLE: compact recall notes. Follow the slide order. Bullets with
%  bold key terms, every formula from the slides, small figures only
%  where the picture is the content. No textbook explanations; add a
%  short line only where a slide is unclear on its own.
% =====================================================================
\lecture{NN}{Lecture Title}{Day, Month DD, 2026}\label{ch:lecNN}

% One-line roadmap (the lecture's own outline slide)
\textbf{Outline.} Topic A $\to$ topic B $\to$ topic C.

% ---------------------------------------------------------------------
\section{Section Title (a divider slide of the deck)}\label{sec:lecNN-a}
\slides{3--9}

\subsection{Sub-topic (a group of slides)}
\slides{3--5}

\begin{itemize}
  \item \textbf{Problem:} softmax normaliser sums over all $|V|$ nodes
    $\Rightarrow$ $\bigO(|V|^2)$ per epoch.
  \item \textbf{Fix -- negative sampling:} normalise against $k$ random
    \term{negative samples} $n_i\sim P_V$ instead of all nodes:
\end{itemize}
\begin{equation}\label{eq:lecNN-negsamp}
  \log\frac{\exp(\vz_u\tp\vz_v)}{\sum_{n\in V}\exp(\vz_u\tp\vz_n)}
  \approx \log\sigma(\vz_u\tp\vz_v)+\sum_{i=1}^{k}\log\sigma(-\vz_u\tp\vz_{n_i})
\end{equation}
\begin{itemize}
  \item Sample negatives $\propto$ degree. $k$: higher $\Rightarrow$ more
    robust, but more bias on negatives; in practice $k=5$--$20$.
\end{itemize}

\begin{definition}[Name]\label{def:lecNN-name}
  Use only for the core formal definitions of the lecture.
\end{definition}

\begin{keyidea}
  Only for a takeaway the slide itself highlights.
\end{keyidea}

\begin{example}
  Only for small worked examples that are on the slides.
\end{example}

\begin{caution}
  Slide erratum (one or two lines): what the slide says $\to$ correct version.
\end{caution}

% Small figure where the picture is the content
\begin{figure}[htbp]
  \centering
  \begin{tikzpicture}[scale=0.9]
    \node[target] (a) at (0,0) {A};
    \node[vertex] (b) at (1.4,0.6) {B};
    \node[vertex] (c) at (1.4,-0.6) {C};
    \draw[edge] (a) -- (b) (a) -- (c) (b) -- (c);
  \end{tikzpicture}
  \caption{Short caption (slide 4).}\label{fig:lecNN-example}
\end{figure}

% ---------------------------------------------------------------------
\bigskip\phantomsection
\addcontentsline{toc}{section}{Lecture NN at a glance}
\begin{summarybox}[Lecture NN at a glance]
  \begin{itemize}
    \item One line per main idea / the key formulas, nothing new.
  \end{itemize}
\end{summarybox}
